\documentclass[12pt]{article}
\usepackage[utf8]{inputenc}
\usepackage{amsmath, amssymb, amsthm}
\usepackage{geometry}
\usepackage{graphicx}
\usepackage{enumitem}
\usepackage{fancyhdr}

\geometry{a4paper, margin=1in}

% Header/Footer Setup
\pagestyle{fancy}
\fancyhf{}
\lhead{Homework 12.2: Vectors}
\rhead{Comprehensive Problem Set}
\cfoot{\thepage}

\title{\textbf{Comprehensive Guide \& Practice: Vectors}\\ \large Based on Homework 12.2 \& Supplementary Concepts}
\author{Math Practice Section}
\date{}

\begin{document}

\maketitle

\section*{Concept Checklist & Study Guide}
Before beginning the problems, review this checklist derived from the provided course materials. Ensure you are comfortable with:

\begin{enumerate}
    \item \textbf{Definitions:} Distinguishing between Scalars (magnitude only) and Vectors (magnitude and direction).
    \item \textbf{Geometric Representation:} Visualizing vectors, the Parallelogram Law, and the Triangle Law.
    \item \textbf{Component Form:} Converting points to vectors ($\text{Terminal} - \text{Initial}$).
    \item \textbf{Magnitude:} Calculating length in 2D ($\sqrt{a^2+b^2}$) and 3D ($\sqrt{a^2+b^2+c^2}$).
    \item \textbf{Vector Arithmetic:} Addition, subtraction, and scalar multiplication.
    \item \textbf{Unit Vectors:} Normalizing vectors ($\mathbf{u} = \mathbf{v}/|\mathbf{v}|$) and standard basis notation ($\mathbf{i}, \mathbf{j}, \mathbf{k}$).
    \item \textbf{Direction \& Angles:} Finding the angle with the positive x-axis and reconstructing vectors from magnitude/direction $(r\cos\theta, r\sin\theta)$.
    \item \textbf{Geometric Applications:} Midpoints, parallel vectors, and geometric proofs using vector algebra.
    \item \textbf{Calculus Connection:} Finding tangent vectors to curves using derivatives ($slope \to \langle 1, m \rangle$).
    \item \textbf{Real-World Applications:} Physics (Force/Velocity) and Finance (Portfolio Vectors).
    \item \textbf{Advanced Concepts (Section 7 of PDF):} Projections, Parametric Equations of Lines, and Linear Combinations.
\end{enumerate}

\newpage

\section*{Part 1: Problem Set}

\subsection*{I. Conceptual Basics (Scalars vs. Vectors)}
\textbf{1.} Classify the following as a scalar or a vector and explain why:
\begin{enumerate}[label=(\alph*)]
    \item The acceleration of a falling object.
    \item The balance of a bank account.
    \item The momentum of a moving train.
\end{enumerate}

\textbf{2.} Given a parallelogram $PQRS$ (labeled counter-clockwise), list two vectors equal to $\vec{PQ}$ and two vectors equal to $\vec{PS}$.

\subsection*{II. Vector Representation and Magnitude}
\textbf{3.} Find the component form of the vector $\mathbf{v}$ with initial point $A(2, -5)$ and terminal point $B(-3, 7)$.

\textbf{4.} Find the vector $\mathbf{w}$ starting at $P(1, 0, -4)$ and ending at $Q(3, 2, 1)$.

\textbf{5.} Calculate the magnitude of the vector $\mathbf{a} = \langle -5, 12 \rangle$.

\textbf{6.} Calculate the magnitude of $\mathbf{b} = 2\mathbf{i} - 4\mathbf{j} + \sqrt{5}\mathbf{k}$.

\subsection*{III. Vector Arithmetic}
\textbf{7.} Given $\mathbf{a} = \langle 3, -1 \rangle$ and $\mathbf{b} = \langle 1, 4 \rangle$, find:
\begin{enumerate}[label=(\alph*)]
    \item $2\mathbf{a} + \mathbf{b}$
    \item $3\mathbf{a} - 2\mathbf{b}$
    \item $|\mathbf{a} - \mathbf{b}|$
\end{enumerate}

\textbf{8.} Given $\mathbf{u} = \mathbf{i} + 3\mathbf{j} - 2\mathbf{k}$ and $\mathbf{v} = 4\mathbf{i} - \mathbf{j}$, find the vector $\mathbf{w} = 2\mathbf{u} + \mathbf{v}$.

\textbf{9.} Find scalars $c_1$ and $c_2$ such that $c_1\langle 3, 2 \rangle + c_2\langle 1, -1 \rangle = \langle 11, 4 \rangle$.

\subsection*{IV. Unit Vectors and Direction}
\textbf{10.} Find a unit vector $\mathbf{u}$ in the direction of $\mathbf{v} = \langle -4, 3 \rangle$.

\textbf{11.} Find a unit vector in the direction of $\mathbf{w} = \langle 2, -1, 2 \rangle$. Express your answer in $\mathbf{i}, \mathbf{j}, \mathbf{k}$ notation.

\textbf{12.} Find a vector of length 5 pointing in the \textit{opposite} direction of $\mathbf{a} = \langle 1, \sqrt{3} \rangle$.

\subsection*{V. Angles and Polar Forms}
\textbf{13.} Find the angle $\theta$ (in degrees) that the vector $\mathbf{v} = \langle 3, 3 \rangle$ makes with the positive x-axis.

\textbf{14.} Find the angle $\theta$ (in degrees, $0^\circ \le \theta < 360^\circ$) that $\mathbf{v} = \langle -1, \sqrt{3} \rangle$ makes with the positive x-axis.

\textbf{15.} A vector $\mathbf{v}$ has magnitude $|\mathbf{v}| = 20$ and makes an angle of $\theta = 150^\circ$ with the positive x-axis. Find the component form of $\mathbf{v}$.

\subsection*{VI. Geometric Connections and Midpoints}
\textbf{16.} (Geometric Proof Setup) In $\triangle ABC$, let $\mathbf{u} = \vec{AB}$ and $\mathbf{v} = \vec{AC}$. Let $M$ be the midpoint of $BC$. Express $\vec{AM}$ in terms of $\mathbf{u}$ and $\mathbf{v}$.

\textbf{17.} In a regular hexagon $ABCDEF$ centered at the origin $O$, let $\vec{OA} = \mathbf{a}$ and $\vec{OB} = \mathbf{b}$. Express $\vec{CD}$ in terms of $\mathbf{a}$ and $\mathbf{b}$.

\subsection*{VII. Calculus Connections (Tangents)}
\textbf{18.} Find a unit vector parallel to the tangent line of the curve $y = x^3 - 4x$ at the point where $x = 2$.

\textbf{19.} Find a direction vector for the tangent line to $y = \sin(x)$ at $x = \pi$.

\subsection*{VIII. Applications (Physics \& Finance)}
\textbf{20.} (Physics) Three forces act on an object: $\mathbf{F}_1 = \langle 10, 0 \rangle$, $\mathbf{F}_2 = \langle 0, -5 \rangle$, and $\mathbf{F}_3 = \langle -2, 8 \rangle$. Find the magnitude of the resultant net force.

\textbf{21.} (Physics/Relative Velocity) An airplane heads North at an airspeed of 400 mph. A wind is blowing from the West (heading East) at 30 mph. Find the ground speed (magnitude of the velocity vector).

\textbf{22.} (Finance - Based on PDF Sec 6) An investor targets a portfolio allocation of 60\% Stocks and 40\% Bonds.
\begin{itemize}
    \item Target vector: $\mathbf{t} = \langle 0.60, 0.40 \rangle$.
    \item Due to market changes, the current weights are $\mathbf{c} = \langle 0.65, 0.35 \rangle$.
    \item If the total portfolio value is $V = \$100,000$, find the "Rebalancing Vector" in dollars (Total Value $\times$ $(\mathbf{t} - \mathbf{c})$).
\end{itemize}

\subsection*{IX. Advanced Variations (PDF Section 7 Topics)}
\textbf{23.} (Parallel Vectors) Determine if $\mathbf{a} = \langle 6, -9, 3 \rangle$ and $\mathbf{b} = \langle -2, 3, -1 \rangle$ are parallel. Explain.

\textbf{24.} (Parametric Line) Find the vector equation of the line passing through $P(2, -1)$ with direction vector $\mathbf{d} = \langle 3, 4 \rangle$.

\textbf{25.} (Parametric Line 3D) Find the parametric equations for the line passing through $A(1, 2, 3)$ and $B(4, 2, -1)$.

\textbf{26.} (Projections) Find the vector projection of $\mathbf{a} = \langle 4, 1 \rangle$ onto $\mathbf{b} = \langle 1, 0 \rangle$. (Hint: $\text{proj}_{\mathbf{b}}\mathbf{a} = \frac{\mathbf{a}\cdot\mathbf{b}}{|\mathbf{b}|^2}\mathbf{b}$).

\textbf{27.} (Linear Span) Can the vector $\mathbf{r} = \langle 5, 7 \rangle$ be written as a linear combination of $\mathbf{i} = \langle 1, 0 \rangle$ and $\mathbf{j} = \langle 0, 1 \rangle$? Write the equation.

\subsection*{X. "Find the Flaw" (Diagnostic)}
\textit{Identify the error in the proposed solution and correct it.}

\textbf{28.} \textbf{Problem:} Find $|\langle -3, -4 \rangle|$.
\textbf{Flawed Solution:} $|\mathbf{v}| = \sqrt{-3^2 + -4^2} = \sqrt{-9 + -16} = \sqrt{-25}$. Undefined.

\textbf{29.} \textbf{Problem:} Find a vector from $A(1,1)$ to $B(4,5)$.
\textbf{Flawed Solution:} $\vec{AB} = A - B = \langle 1-4, 1-5 \rangle = \langle -3, -4 \rangle$.

\textbf{30.} \textbf{Problem:} Calculate the magnitude of $5\mathbf{v}$ where $|\mathbf{v}|=2$.
\textbf{Flawed Solution:} $|5\mathbf{v}| = 5 + |\mathbf{v}| = 5 + 2 = 7$.

\textbf{31.} \textbf{Problem:} Find the angle of $\mathbf{v} = \langle -1, -1 \rangle$.
\textbf{Flawed Solution:} $\tan \theta = \frac{-1}{-1} = 1$. $\theta = \arctan(1) = 45^\circ$.

\textbf{32.} \textbf{Problem:} Normalize $\mathbf{v} = \langle 3, 4 \rangle$.
\textbf{Flawed Solution:} $\mathbf{u} = |\mathbf{v}| \cdot \mathbf{v} = 5 \langle 3, 4 \rangle = \langle 15, 20 \rangle$.

\newpage

\section*{Part 2: Detailed Solutions}

\textbf{1.}
(a) \textbf{Vector.} Acceleration has magnitude ($9.8 m/s^2$) and direction (down).
(b) \textbf{Scalar.} Money has amount but no spatial direction.
(c) \textbf{Vector.} Momentum is mass $\times$ velocity (velocity is a vector).

\textbf{2.} Opposite sides of a parallelogram are equal vectors.
$\vec{PQ} = \vec{SR}$ (same length, parallel, same direction).
$\vec{PS} = \vec{QR}$.

\textbf{3.} $\mathbf{v} = \text{Terminal} - \text{Initial} = \langle -3 - 2, 7 - (-5) \rangle = \langle -5, 12 \rangle$.

\textbf{4.} $\mathbf{w} = \langle 3-1, 2-0, 1-(-4) \rangle = \langle 2, 2, 5 \rangle$.

\textbf{5.} $|\mathbf{a}| = \sqrt{(-5)^2 + 12^2} = \sqrt{25 + 144} = \sqrt{169} = 13$.

\textbf{6.} $|\mathbf{b}| = \sqrt{2^2 + (-4)^2 + (\sqrt{5})^2} = \sqrt{4 + 16 + 5} = \sqrt{25} = 5$.

\textbf{7.}
(a) $2\langle 3, -1 \rangle + \langle 1, 4 \rangle = \langle 6, -2 \rangle + \langle 1, 4 \rangle = \langle 7, 2 \rangle$.
(b) $3\langle 3, -1 \rangle - 2\langle 1, 4 \rangle = \langle 9, -3 \rangle - \langle 2, 8 \rangle = \langle 7, -11 \rangle$.
(c) $\mathbf{a}-\mathbf{b} = \langle 3-1, -1-4 \rangle = \langle 2, -5 \rangle$. Magnitude: $\sqrt{2^2+(-5)^2} = \sqrt{29}$.

\textbf{8.} $\mathbf{u} = \langle 1, 3, -2 \rangle$, $\mathbf{v} = \langle 4, -1, 0 \rangle$.
$2\mathbf{u} + \mathbf{v} = \langle 2, 6, -4 \rangle + \langle 4, -1, 0 \rangle = \langle 6, 5, -4 \rangle$.

\textbf{9.} System of equations:
$3c_1 + 1c_2 = 11$ (x-components)
$2c_1 - 1c_2 = 4$ (y-components)
Add equations: $5c_1 = 15 \implies c_1 = 3$.
Sub into first: $3(3) + c_2 = 11 \implies 9 + c_2 = 11 \implies c_2 = 2$.
Solution: $c_1=3, c_2=2$.

\textbf{10.} $|\mathbf{v}| = \sqrt{16+9} = 5$. Unit vector $\mathbf{u} = \frac{1}{5}\langle -4, 3 \rangle = \langle -0.8, 0.6 \rangle$.

\textbf{11.} $|\mathbf{w}| = \sqrt{4+1+4} = \sqrt{9} = 3$. Unit vector = $\frac{2}{3}\mathbf{i} - \frac{1}{3}\mathbf{j} + \frac{2}{3}\mathbf{k}$.

\textbf{12.} $|\mathbf{a}| = \sqrt{1+3} = 2$. Unit vector direction $\mathbf{u}_a = \langle 1/2, \sqrt{3}/2 \rangle$.
Opposite direction $= -\mathbf{u}_a$. Length 5 means multiply by 5.
Result: $-5 \langle 1/2, \sqrt{3}/2 \rangle = \langle -2.5, -2.5\sqrt{3} \rangle$.

\textbf{13.} Quadrant I. $\tan \theta = 3/3 = 1$. $\theta = 45^\circ$.

\textbf{14.} $x=-1, y=\sqrt{3}$. Quadrant II. Reference angle $\alpha = \arctan(|\sqrt{3}/-1|) = 60^\circ$.
$\theta = 180^\circ - 60^\circ = 120^\circ$.

\textbf{15.} $x = 20 \cos(150^\circ) = 20(-\sqrt{3}/2) = -10\sqrt{3}$.
$y = 20 \sin(150^\circ) = 20(1/2) = 10$.
$\mathbf{v} = \langle -10\sqrt{3}, 10 \rangle$.

\textbf{16.} $\vec{AM} = \vec{AB} + \vec{BM}$. Since $M$ is midpoint, $\vec{BM} = \frac{1}{2}\vec{BC}$.
Also $\vec{BC} = \vec{AC} - \vec{AB} = \mathbf{v} - \mathbf{u}$.
$\vec{AM} = \mathbf{u} + \frac{1}{2}(\mathbf{v} - \mathbf{u}) = \mathbf{u} - \frac{1}{2}\mathbf{u} + \frac{1}{2}\mathbf{v} = \frac{1}{2}(\mathbf{u} + \mathbf{v})$.

\textbf{17.} In a regular hexagon centered at origin, $\vec{CD}$ is parallel and equal to $\vec{AF}$ and $\vec{BO}$ (pointing inward).
Actually, simpler geometry: $\vec{CD} = \vec{CO} + \vec{OD}$.
Note $\vec{C}$ is opposite $\vec{F}$, $\vec{D}$ is opposite $\vec{A}$. This depends on labeling.
Standard alternative: $\vec{CD}$ is parallel to $\vec{BO} = -\mathbf{b}$.
If regular, $\vec{CD} = \vec{BO} = -\mathbf{b}$.

\textbf{18.} $y = x^3 - 4x \implies y' = 3x^2 - 4$.
At $x=2$: $y' = 3(4) - 4 = 8$. Slope $m=8$.
Vector direction $\mathbf{d} = \langle 1, 8 \rangle$. Magnitude $\sqrt{1+64} = \sqrt{65}$.
Unit vector: $\langle \frac{1}{\sqrt{65}}, \frac{8}{\sqrt{65}} \rangle$.

\textbf{19.} $y' = \cos(x)$. At $x=\pi$, $y' = -1$.
Slope vector $\mathbf{v} = \langle 1, -1 \rangle$.

\textbf{20.} Net Force $\mathbf{F}_{net} = \mathbf{F}_1 + \mathbf{F}_2 + \mathbf{F}_3$.
x-comp: $10 + 0 - 2 = 8$.
y-comp: $0 - 5 + 8 = 3$.
$\mathbf{F}_{net} = \langle 8, 3 \rangle$. Magnitude $= \sqrt{64+9} = \sqrt{73} \approx 8.54$ N.

\textbf{21.} Plane $\mathbf{p} = \langle 0, 400 \rangle$ (North). Wind $\mathbf{w} = \langle 30, 0 \rangle$ (East).
Ground velocity $\mathbf{v} = \langle 30, 400 \rangle$.
Speed $= |\mathbf{v}| = \sqrt{30^2 + 400^2} = \sqrt{900 + 160000} = \sqrt{160900} \approx 401.1$ mph.

\textbf{22.} Diff $\mathbf{d} = \mathbf{t} - \mathbf{c} = \langle 0.60 - 0.65, 0.40 - 0.35 \rangle = \langle -0.05, 0.05 \rangle$.
Rebalance Dollars $= 100,000 \times \langle -0.05, 0.05 \rangle = \langle -5000, 5000 \rangle$.
(Sell \$5,000 Stock, Buy \$5,000 Bonds).

\textbf{23.} Check ratios: $-2/6 = -1/3$. $3/-9 = -1/3$. $-1/3 = -1/3$.
Since ratios are constant, $\mathbf{b} = -\frac{1}{3}\mathbf{a}$. Yes, they are parallel.

\textbf{24.} Vector equation: $\mathbf{r}(t) = \mathbf{r}_0 + t\mathbf{d}$.
$\langle x, y \rangle = \langle 2, -1 \rangle + t\langle 3, 4 \rangle$.

\textbf{25.} Direction $\mathbf{d} = B - A = \langle 3, 0, -4 \rangle$.
Line: $x = 1 + 3t$, $y = 2 + 0t$, $z = 3 - 4t$.

\textbf{26.} Dot product $\mathbf{a} \cdot \mathbf{b} = 4(1) + 1(0) = 4$.
$|\mathbf{b}|^2 = 1^2 + 0 = 1$.
Proj $= \frac{4}{1} \langle 1, 0 \rangle = \langle 4, 0 \rangle$.

\textbf{27.} Yes. $c_1\langle 1,0 \rangle + c_2\langle 0,1 \rangle = \langle c_1, c_2 \rangle$.
Set $c_1=5, c_2=7$. Equation: $\mathbf{r} = 5\mathbf{i} + 7\mathbf{j}$.

\textbf{28.} \textbf{Error:} Squaring negative numbers yields positives. $(-3)^2 = 9$.
\textbf{Correction:} $\sqrt{9 + 16} = 5$.

\textbf{29.} \textbf{Error:} Subtracted Start - End. Must be End - Start.
\textbf{Correction:} $\vec{AB} = \langle 4-1, 5-1 \rangle = \langle 3, 4 \rangle$.

\textbf{30.} \textbf{Error:} Scalar multiplication scales magnitude, it doesn't add.
\textbf{Correction:} $|5\mathbf{v}| = |5| \cdot |\mathbf{v}| = 5 \times 2 = 10$.

\textbf{31.} \textbf{Error:} Quadrant issue. $\langle -1, -1 \rangle$ is in Q3. Calculator gives Q1 ($45^\circ$).
\textbf{Correction:} $45^\circ + 180^\circ = 225^\circ$.

\textbf{32.} \textbf{Error:} Multiplied by magnitude instead of dividing.
\textbf{Correction:} Magnitude is 5. $\mathbf{u} = \langle 3/5, 4/5 \rangle$.

\newpage

\section*{Part 3: Concept Coverage Map}
Use this table to double-check your understanding of specific topics found in the homework files.

\begin{center}
\begin{tabular}{|l|l|}
\hline
\textbf{Concept / Topic} & \textbf{Problem Numbers} \\
\hline
Scalar vs. Vector Identification & 1 \\
\hline
Geometric Vector Properties & 2, 17 \\
\hline
Vector from Two Points & 3, 4, 29 \\
\hline
Magnitude (Distance) & 5, 6, 28, 30 \\
\hline
Vector Arithmetic (Add/Sub/Scale) & 7, 8, 20 \\
\hline
Linear Combinations / Solving & 9, 27 \\
\hline
Unit Vectors & 10, 11, 32 \\
\hline
Angle / Polar Forms & 13, 14, 15, 31 \\
\hline
Geometric Proofs / Midpoints & 16 \\
\hline
Tangent Vectors (Calculus) & 18, 19 \\
\hline
Physics Applications & 20, 21 \\
\hline
Finance Applications & 22 \\
\hline
Parallel Vectors & 12, 23 \\
\hline
Lines and Planes (Parametric) & 24, 25 \\
\hline
Projections (Dot Product intro) & 26 \\
\hline
\end{tabular}
\end{center}

\end{document}